\input{preamble}
\begin{document}

\title{Bridgeland Stability: the General
Theory}
\maketitle

Minicourse by Cristian Martinez (Universidad
de los Andes, Bogot\'a, Colombia), CIMPA
school
Florian\'opolis 2025.

Notes at
\href{%
http://github.com/danimalabares/cimpa-floripa
}{%
github.com/danimalabares/cimpa-floripa
}

\bigskip\noindent

{\bf Abstract.} Bridgeland stability is a
powerful tool for extracting geometry from
homological algebra. In particular, it gives
a framework for studying moduli spaces of
objects in a triangulated category, such as
the derived category of an algebraic variety.
The subject was born as a mathematical
interpretation of work in string theory, but
has since impacted many areas, including
classical algebraic geometry, derived
categories of coherent sheaves, enumerative
geometry, homological mirror symmetry, and
symplectic geometry. The goal of this course
is to develop the foundations of Bridgeland
stability, covering the following topics: 1)
The theory of t-structures on triangulated
categories, including examples via tilting.
2) The defniition of stability conditions and
the stability manifold, as well as
Bridgeland's deformation theorem. 3)
Constructions of stability conditions. 4)
Moduli spaces of stable objects.

\bigskip\noindent
\tableofcontents
\bigskip\noindent

\noindent
{\bf Plan.}
\begin{enumerate}
\item Lecture 1: What is a stability
condition.
\item Lecture 2: Main constructions in
dimensions 2 and 3.
\item Lecture 3: Reaching the Geiseker
chamber.
\item Lecture 4: Some applications on
structures.
\end{enumerate}

\section{Recall}
\label{section-recall}

Recall that for $\mu_H$-stability we have the
following property.

\begin{itemize}
\item (Schur's lemma.) If $A,B \in
\text{Coh}$ are $\mu_H$ semistable with
$\mu_H(A)>\mu_H(B)$ then $\Hom(A,B)=0$.
Exercise.
\item (Harder-Narasimhan filtrations.) $E \in
\text{Coh}(X)$ then is a
filtration
$$
\xymatrix{
0\ar[r]&E_0\ar[rr]^{i_0}&&E_1\ar[dl]\ar[r]&
\cdots\ar[r]
&E_{k-1}\ar[rr]^{i_k}&&E_k=E\ar[dl]\\
&&F_1& & & & F_k
}
$$
\iffalse$$
\xymatrix{
0\ar[r]&  E_0
\ar@{^{(}->}[r]&E_1\ar[d]\ar[r]&  \cdots
\ar@{^{(}->}[r]& E_k\ar[d]=E\\ & & F_1 & &
F_k } $$\fi with $E_0$ is the torsion
subsheaf of $ E$, $F_j$ are
$\mu_H$-semistable with
$$
\mu_H(F_1)>\mu_H(F_2)>\ldots>\mu_H(F_k)
$$
\end{itemize}

Exercise: (1) $\implies$ (2).

These two results basically say that
$\text{Coh}(X)$ is generated by
$\mu_H$-semistable sheaves (via extensions)
and torsion sheaves.
$$
0\to E_{j-1}\hookrightarrow E_j\to F_j\to_0
$$
$E_0:=F_0$.

Can we do this on $D^b(X)$?

We look for a subcategory with a fixed slope
to generate the whole category\ldots{}

\medskip

The following definition is intended to
``impose'' Schur's lemma on our objects:

\begin{definition}
\label{definition-slicing}

A {\it slicing} $\mathcal{P}$ of $D^b(X)$
consists of full subcategories
$\mathcal{P}(\phi)$ for every $\phi \in
\mathbb{R}$ such that
\begin{enumerate}
\item
$\mathcal{P}(\phi)[1]=\mathcal{P}(\phi+1)$ /
\item If $\phi_1>\phi_2$ then
$\Hom(\mathcal{P}(\phi_1),\mathcal{P}(\phi_2)
)
=0$.
\item (HN property.) For every $E \in D^b(X)$
there is a diagram (filtration)
$$
\phi_1>\phi_2>\cdots>\phi_k
$$
$$
\xymatrix{
0=E_0\ar[rr]^{i_0}&&E_1\ar[dl]\ar[r]&
\cdots\ar[r]
&E_{k-1}\ar[rr]^{i_k}&&E_k=E\ar[dl]\\
&F_1& & & & F_k
}
$$
where $F_j=\text{Cone}(i_j) \in
\mathcal{P}(d_j)$. {\bf Exercise:} this is
unique because of (2).

\item  $F \in Q(\phi)$ is called {\it
semistable} of phase $\phi$. The simple
objects
(those which have no subobjects) in
$\mathcal{P}(\phi)$ are called {\it stable of
phase} $\phi$.

\item  $\phi_1:=\phi^+(E)$,
$\phi_k:=\phi^-(E)$.
\end{enumerate}
\end{definition}

The following definitions is related to the
former, and to Schur's lemma:

\begin{definition}
\label{bsgt-label-001}
The {\it heart of a bounded $t$-structure} on
$D^b(X)$ is a full additive subcategory
$\mathcal{A} \subset D^b(X)$ such that
\begin{enumerate}
\item For $i>j$, $\Hom(A[i],B[j])=0$ for $A,B
\in \mathcal{A}$.
\item For every $E \in D^b(X)$ there are
integers $k_1>k_2>\ldots>k_m$ and a
filtration
\begin{equation}
\label{equation-diagram-abelian-category}
\xymatrix{
0=E_0\ar[rr]^{i_0}&&E_1\ar[dl]\ar[r]&
\cdots\ar[r]
&E_{k-1}\ar[rr]^{i_k}&&E_k=E\ar[dl]\\
&A_1[k_1]& & & & A_k[k_m]
}
\end{equation}
exact triangles, with $A_j \in \mathcal{A}$.
\end{enumerate}
\end{definition}

\begin{exercise}
\label{exercise-A-is-abelian}
$\mathcal{A}$ is abelian.
\end{exercise}

\begin{exercise}
\label{exercise-}
If $\mathcal{P}$ is a slicing of $D^b(X)$
then $\mathcal{P}(0,]$ (the extension closure
of $\{\mathcal{P}(\phi):0<\phi \leq 1\}$) is
the heart of a bounded $t$-structure.
\end{exercise}

\begin{example}
\label{bsgt-label-002}
$\text{Coh}(X)\subset D^b(X)$ is the heart of
a bounded structure.

$$
\xymatrix{
\sigma_{\leq i}(E^\bullet):\ar[d]&
\cdots\ar[r]& E^{j-2}\ar[d]_=\ar[r]&
E^{j-1}\ar[d]_=\ar[r]^{\delta}& \Ker
\delta_j\ar[d]\ar[r]&  0 \ar[d]\ar[r]
&  \cdots\\
\sigma_{\leq j+1}(E^\bullet):\ar[d]&  \cdots
\ar[r]& E^{j-2}\ar[d]\ar[r]
& E^{j-1}\ar[d]\ar[r]^{d_j}&E^j\ar[d]\ar[r]&
\Ker(\delta_{j+1})\ar[d]\\
&& 0& 0&E_2k/\Ker \delta_j& \Ker \delta_{j+1}
}
$$
\medskip
$$
\xymatrix{
\cdots\ar[r]&\sigma_{\leq
j}(E^\bullet)\ar[rr]& &\sigma_{\leq
j+1}(E^\bullet)\ar[dl]\ar[r]&  \cdots\\
&  &  H^{i+j}(E^\bullet)[-j-1]
}
$$
\end{example}

\begin{definition}
\label{definition-cohomologies}
The objects $A_j$ appearing in Eq.
\ref{equation-diagram-abelian-category} are
called the {\it cohomologies} of $E$ with
respect to the $t$-structure $\mathcal{A}$,
$A_j:=\mathcal{H}_{\mathcal{A}}^{-k_j}(E)$.
\end{definition}

\section{Constructing new hearts}
\label{section-constructing-new-heats}

How can we construct other $t$-structures?
Suppose that $\mathcal{A}$ is an abelian
category. A {\it torsion pair} on
$\mathcal{A}$ consists of two full
subcategories $(\mathcal{T}, \mathcal{F})$
such that
\begin{enumerate}
\item $\Hom(T,F)=0$ for all $T \in
\mathcal{T}$ and $F \in \mathcal{F}$.
\item For every $A \in \mathcal{A}$ there is
an exact sequence
$$
\xymatrix{
0\ar[r]&T\ar[r]&A\ar[r]&F\ar[r]&0
}
$$
with $T \in \mathcal{T}$ and $F \in
\mathcal{F}$.
\end{enumerate}

Again: with the first property we may prove
that the exact sequence in the second
property is unique (Exercise).

\begin{exercise}
\label{exercise-torsion-pair}
$\mathcal{A}=\text{Coh}(X)$, $\mathcal{T}=$
torsion sheaves, $\mathcal{F}=$ torsion free
sheaves.
\end{exercise}

\begin{proposition}
\label{proposition-}
If $\mathcal{A}$ is a hear of a bounded
$t$-structure on $D^b(X)$ and
$(\mathcal{T},\mathcal{F})$ is a torsion pair
on $\mathcal{A}$, then
\begin{align*}
\mathcal{A}^\sharp&=\left<\mathcal{F}[1],
\mathcal{T}\right>\\
&=\left\{E\in
D^b(X):\substack{\mathcal{H}_{\mathcal{A}}
^j(E)=0\\
\text{for $j \neq -1,0$} \\
\mathcal{H}_{\mathcal{A}}^{-1}(E)=\mathcal{F}
\\
\mathcal{H}_{\mathcal{A}}^\bullet(E) in
\mathcal{T} }\right\}
\end{align*}
which is called the {\it tilted heart
category}.
\end{proposition}

\noindent
{\bf Note.} We can think of those as
complexes $E ^\bullet$ that fit in an exact
triangle
$$
\xymatrix{
0\ar[r]&F[1]\ar[r]&E^\bullet\ar[r]&T\ar[r]&0
}
$$
(which a prori is just an exact triangle, but
it is an exact sequence in our case).

The first proof of this fact was done by
Bridgeland [confirm]:

\begin{proof}[Idea of proof]
[Picture.]
\end{proof}

We show an example to understand how the
proof works.

\begin{example}
\label{example-}
Suppose that the $\mathcal{A}$-filtration of
an object $E$ is of the form
$$
\xymatrix{
0\ar[r]&F_1[1]\ar[r]&E\ar[r]&F_2\ar[r]&0
}
$$

Let's comute the cohomologies:
$$
\xymatrix{
0=E_0\ar[rr]^{i_0}&&E_1\ar[dl]\ar[rr]^{i_2}&&
E_2=E\ar[dl]\\
&F_1[1]&  & F_2
}
$$

$F_1,F_2 \in \mathcal{A}$. So, the
cohomologies, by definition are
$\mathcal{H}_{\mathcal{A}}^{-1}(E)=F_1$,
$\mathcal{H}_{\mathcal{A}}^0(E)=F_2$.

Now consider the following diagram:
$$
\xymatrix{
&  F[1]\ar[d]&  &  T\ar[d]\\
0 \ar[r]&  F_1[1]\ar[d]\ar[r]&  E
\ar[r]\ar[rd]_f&  F_2\ar[d]\ar[r]& 0\\
& F[1]& & F
}
$$
so wer consider
$$
\xymatrix{
& & F[1]\ar[dl]\ar[d]\\
0\ar[r]&C=\text{Cone}(f)[-1]\ar[r]&E\ar[r]^f&
F\ar[r]&0
}
$$
\end{example}

Recall that in a triangulated category,
$$
\xymatrix{
X\ar[d]\ar[r]&  0 \ar[d]\ar[r]&
X[1]\ar@{.>}[d]\ar[r]& X[1]\ar[d]\\
Y\ar[r]& Z\ar[r]& X^1\ar[r]&Y[1]
}
$$
where the leftmost square commutes.

$$
\xymatrix{
0\ar[rr]&&F[1]\ar[dl]\ar[rr]^{g}&  &  C
\ar[dl]\ar[rr]&  &  E\ar[dl]\\
&  \underbrace{F[1]}_{\mathcal{A}^\sharp[1]}&
&
\underbrace{\text{Cone}(g)}_{\mathcal{A}
^\sharp}&
&
\underbrace{F}_{\mathcal{A}^\sharp[-1]}
}
$$

\section{Stability conditions}
\label{section-stability-conditions}

\begin{definition}[Stability condition,
definition 1]
\label{definition-stability-condition-1}
A {\it stability condition} is a pair
$\sigma=(\mathcal{P},Z)$ where $\mathcal{P}$
is a slicing of $D^b(X)$, and a morphism $Z$
called the  {\it central change} of the form
$$
\xymatrix{
Z:K_0(D^b(X))\ar[rr]\ar@{->>}[dr]_\nu&  &
\mathbb{C}\\
&  \Gamma \ar[ur]_z
}
$$
where $\Gamma$ is a finite rank lattice,
satisfying that
\begin{enumerate}
\item $Z(E) \in \mathcal{E}_{>0}e^{i\pi
\phi}$ for any nonzero object $E \in
\mathcal{P}(\phi)$.
\item (Support property.)
For a fixed norm $\|\cdot\|$ on $\Gamma
\otimes \mathbb{R}$ we have
$$
C_{\sigma}=\text{inf}\left\{
\frac{|Z(E)|}{\|\nu(E)\|}:\substack{
\text{$E$ non-zero} \\ E \in
\mathcal{P}(\phi),\phi \in \mathbb{R}}
\right\} >0
$$
\end{enumerate}
\end{definition}

\begin{definition}[Stability condition,
definition 2]
\label{definition-stability-condition-2}
A {\it stability condition} is a pair
$\sigma=(\mathcal{A},Z)$, where $\mathcal{A}$
is the heart of a bounded $t$-structure on
$D^b(X)$, and a morphism $Z$ of the form
$$
\xymatrix{
Z:K_0(D^b(X))\ar[rr]\ar@{->>}[dr]_\nu&  &
\mathbb{C}\\
&  \Gamma \ar[ur]_Z
}
$$
where $\Gamma$ is a finite rank lattice,
satisfying that
\begin{enumerate}
\item $Z(E) \in \mathbb{R}_{>0}e^{i\pi \phi}$
for $E \in \mathcal{A}$,
$$
\left[
\frac{\text{Im}Z(E)=0}{\text{Im}Z(E)\geq
0}\implies
\text{Re}Z(E)<0 \right]
$$
\item If
$\mu_Z=-\frac{\text{Re}Z}{\text{Im}Z}$ then
objects in $\mathcal{A}$
have the HN property w.r.t. $\mu_Z$ and we
have the support property.
\end{enumerate}
\end{definition}

\begin{exercise}
\label{exercise-definitions-are-equivalent}
Prove that the two definitions of stability
are equivalent.
\begin{itemize}
\item Prove Schur's Lemma for
$\mu_Z$-stability.
\item $\mathcal{P}(\phi)=\mu_Z$-semistable
objects of phase $\phi$, i.e., $Z(E)
\in \mathbb{R}_{>0}e^{i\pi\phi}$ for $\phi
\in (0,1]$ and set
$\mathcal{P}(\phi+1)=\mathcal{P}(\phi)[1]$.
\item Take $\mathcal{A}(0,1]=$generated by
$\mathcal{P}(\phi)$ with
$0<\phi\leq 1$.
\end{itemize}
\end{exercise}

\begin{exercise}
\label{bsgt-label-003}
If $\dim X=1$, $\mathcal{A}=\text{Coh}(X)$,
$Z= - \text{deg}+i \text{rank}$,
$v=(\text{deg},\text{rank})$.
\end{exercise}

\medskip\noindent
{\bf Rephrasing of support property.} There
exists a quadratic form $Q$ on $\Gamma
\otimes \mathbb{R}$ such that $Q(v(E))\geq 0$
for $E$ $\mu_Z$-semistable, If $Z(V)=0$ then
$Q(V)<0$ for nonzero $V$. This is like
imposing Bogomolov inequality. (See
moduli-spaces-of-sheaves.)

\medskip\noindent
Let $\text{Stab}^\Gamma(X)$ be the set of
stability conditions with central charge
factoring through
$v:K_0(\mathcal{A})\to\Gamma$.

\begin{theorem}[Bridgeland]
\label{theorem-Bridgeland}
$\text{Stab}^\Gamma(X)$ is a complex
manifold. Moreover,
\begin{align*}
\text{Stab}^\Gamma(X) &\longrightarrow
\Hom(\Gamma,\mathbb{C}) \\
\sigma=(\mathcal{A},Z) &\longmapsto Z
\end{align*}
is a local homomorphism.
\end{theorem}

This is known as Bridgeland deformation
result. In fact, it is a consequence of the
existence of the quadratic form [reference?].

For every object in $\mathcal{A}$ where
$\sigma=(\mathcal{A},Z)$ is a stability
condition there is a HN filtration
$$
\xymatrix{
0\ar[r]&E_0\ar[rr]^{i_0}&&E_1\ar[dl]\ar[r]&
\cdots\ar[r]
&E_{m-1}\ar[rr]^{i_m}&&E_m=E\ar[dl]\\
&&F_1& & & & F_m
}
$$
$F_1 \in \mathcal{P}(\phi_1)$,
$\phi_1:=\phi^+_\sigma(E)$, $F_m \in
\mathcal{P}(\phi_m)$,
$\phi_m:=\phi_\sigma^-(E)$.

The topology on $\text{Stab}^\Gamma(X)$ is
the coarsest topology that makes the
functions
$$
\sigma \mapsto \phi^+_\sigma,\qquad \sigma
\mapsto \phi^-_\sigma, \qquad
\sigma \mapsto Z
$$
[continuous].

\medskip\noindent
Now let $\dim X=2$. Let $H,B \in
\text{NS}(X)_\mathbb{R}$, $H$ ample.
$$
\text{ch}^B(E)=e^{-B}\text{ch}(E)=
\text{``$\text{ch}(E\otimes
(-B))$''}
$$
When computing the Chern characters we obtain
\begin{align*}
\text{ch}_0^B = \text{ch}_0,&\qquad
\text{ch}_1^B=\text{ch}_1-B\text{ch}_0\\
\text{ch}_2^B&=\text{ch}_2-B\text{ch}_1+
\frac{B^2}{2}\text{ch}_0
\end{align*}

\medskip\noindent
{\bf $B$-Twisted Mumford slope.}
$$
\mu_{B,H}=
\begin{cases}
\frac{H \text{ch}_1^B}{H^2 \text{ch}_0}\qquad
&\text{if }\text{ch}_0\neq 0 \\
+\infty\qquad &\text{if }\text{ch}_0=0
\end{cases}
$$
\begin{align*}
\mu_{B,H}(E)&=\frac{H(\text{ch}_1(E)
-B\text{ch}_0(E)}{H^2\text{ch}_0(E)}\\
&=\mu_H(E)-BH
\end{align*}
$$
\mathcal{T}_{B,H}=\{E \in
\text{Coh}(X):\mu_H\text{-HN factors have }
\mu_{B,H}>0\}
$$
$$
\mathcal{F}_{B,H}=\{E \in \text{Coh}(X):
\mu_H\text{-HN factors have
}\mu_{B,H}\leq 0\}
$$
\begin{itemize}
\item $\Hom(T,F)=0$ for $T \in
\mathcal{T}_{B,H}$, $F \in \mathcal{F}_{B,H}$
(Schur's lemma).
\item If $\mathcal{E}\in\text{Coh}(X)$, take
its HN-filtration
$$
\xymatrix{
0\ar[r]&E_0\ar[rr]^{i_0}&&E_1\ar[dl]\ar[r]&
\cdots\ar[r]
&E_{m-1}\ar[rr]^{i_m}&&E_m=\mathcal{E}\ar[dl]
\\
&&F_1& & & & F_m
}
$$
$$
\mu_{B,H}(E_0)>\mu_{B,H}(F_1)>\ldots>\mu_{B,
H}
(F_m)
$$
Let $j$ such that $\mu_{B,H}(F_j)>0$ and
$\mu_{B,H}\leq 0$.
$$
0 \to \underbrace{E_j}_{\mathcal{T}_{B,H}}\to
\mathcal{E}
\to \underbrace{\mathcal{E}/E_j}_{\in
\mathcal{F}_{B,H}}\to 0
$$
\end{itemize}

Set a tilted heart
$\text{Coh}^{B,H}(X)=\left<\mathcal{F}_{B,H}
[1],
\mathcal{T}_{B,H}\right>$.

Generators:
\begin{itemize}
\item $T$ torsion sheaves $\to
H\text{ch}_1^B(T)\geq 0$.
\item $\mathcal{E}$ $\mu_{B,H}$-semistable
sheaves with $\mu_{B,H}>0 \to
H\text{ch}_1^B (\mathcal{E})>0$.
\item $F[1]$,  $F$ is $\mu_{B,H}$-semistable
with $\mu_{B,H}\leq 0$.
\end{itemize}
Then
$$
H\text{ch}_1^B (F)\leq 0,\qquad
H\text{ch}_1^B (F[1])\geq 0
$$
since shifting by 1 changes the sign of the
Chern character. [Comment about imaginary
part of stability condition.]

Is there a central charge $Z$ such that
$\sigma_{B,H}=(\text{Coh}^{B.H}(X),Z)$ is a
stability condition?
$$
Z_{B,H}=\text{Re}Z_{B,H}+i H\text{ch}_1^B
$$
We need: $H\text{ch}_1^B(E)=0$ for some $E
\in \text{Coh}^{B,H}(X)$ then
$\text{Re}Z_{B,H}(E)<0$.

Ingredients:
\begin{enumerate}
\item (Hodge index theorem.) $(DH)^2 \geq D^2
H^2$.
\item (Bogomolov inequality.) If $E$ is
$\mu_{B,H}$-semistable then
\begin{align*}
\Delta(E)&=[\text{ch}_1(E)]^2-2\text{ch}_0(E)
-\text{ch}_2(E)\geq
0\\
&=\text{ch}_1^B(E)^2-2\text{ch}_0^B\text{ch}
_2^B(E)\geq
0
\end{align*}
We obtain the {\it BH-discriminant}
$$
\Delta_{B,H}(E):=(H\text{ch}_1^B(E))
^2-2\text{ch}_0(E)H^2\text{ch}_2^B(C)
\geq H^2\Delta(E)\geq 0
$$
If $E$ is $\mu_H$-semistable and
$\text{Im}Z_{B,H}(E)=0$, $H\text{ch}_1^B(E)=0
\implies \text{ch}_2^B(E)\leq 0$,
$\text{ch}_2^B(E)-a\text{ch}_0(E)<0$
for any $a>0$.
\end{enumerate}

\begin{theorem}[Bridgeland, Areni-Bestum]
\label{theorem-Bridgeland-Areni-Bersum}
$$
Z_{B,H}=-(\text{ch}_2^B-a\text{ch}_0)+
iH\text{ch}_1^B
$$
is the central charge of a stability
condition on $\text{Coh}^{B,H}(X)$ for any
$a>0$.
\end{theorem}

\begin{theorem}
\label{theorem-BH-discriminant}
$\Delta_{B,H}\geq 0$ for any
$v_{B,H,a}$-semistable object, where
$v_{B,H,a}=\frac{\text{ch}_2^B-a\text{ch}_0}
{H\text{ch}_1^B}$.
\end{theorem}

\medskip\noindent
For $n:=\dim X\geq 2$,
$$
\mu_{B,H}=
\begin{cases}
\frac{H^{n-1}\text{ch}_1^B}{\text{ch}_0}
\qquad
&\text{if }\text{ch}_0 \neq 0 \\
+\infty\qquad &\text{if }\text{ch}_0=0
\end{cases}
$$
$\text{Coh}^{B,H}(X)$, $\Delta_{B,H}(E)=\dim
n$ analogy. $Z_{B,H,a}$ does not satisfy the
positivity property of a central charge since
$$
Z_{B,H,a}=-(H^{n-2}\text{ch}_2^B-a\text{ch}
_0+
iH^{n-1}\text{ch}_1^B)
$$
[skyscraper sheaf argument]

\begin{exercise}
\label{exercise-E-is-a-sheaf}
Let $E \in \text{Coh}^{B,H}(X)$. If
$\text{ch}_0(E)>0$ and $E$ is
$\nu_{B,H,a}$-semistable for all $a\gg 0$,
then $E$ is a sheaf. Moreover, $E$ is
$\mu_{B,H}$-semistable.
\end{exercise}

\begin{proposition}[Lo,---]
\label{proposition-}
If $a>\mu_{B,H}(E)\Delta_{B,H}(E)$, then if
$E$ is {\it twisted} semistable sheaf then
$E$ is $\nu_{B,H,a}$-semistable.
\end{proposition}

\medskip\noindent
If $\dim X=3$ consider the subcategories of
$\text{Coh}^{B,H}(X)$
$$
T_{B,H,a}=\{E \in
\text{Coh}^{B,H}(X):\text{HN factors with }
\nu_{B,H,a}>0\}
$$
$$
\mathcal{F}_{B,H,a}=\{E \in
\text{Coh}^{B,H}(X):\text{HN factors with }
\nu_{B,H,a}\leq 0\}
$$
$$
\mathcal{A}=\left<\mathcal{F}_{B,H,a}[1],
\mathcal{T}_{B,H,a}\right>
$$
supports a stability condition.

\section{GBG inequality conjecture}
\label{section-GBG-inequality-conjecture}

Here's the construction so far:
$$
\xymatrixrowsep{1.4em}
\xymatrix{
\textit{Coh}(X)\ar[r]^{\substack{\text{first}
\\
\text{tilt}}}
&
\textit{Coh}^{B,H}(Z)\ar[r]
^{\substack{\text{second}
\\ \text{tilt}}}
&  \mathcal{A}^{B,H,\alpha}\\
\underbrace{\mu_{B,H}}_{\text{Mumford
stability}}
&
\underbrace{\nu_{B,H,\alpha}}_{\text{Tilt}}
&  \underbrace{\lambda_{B,H,\alpha,s}}
_{\substack{\text{Bridgeland} \\
\text{stability}\\\text{when possible}}}\\
\Delta \geq 0
&  \Delta_{B,H}\geq 0
&\substack{\text{Conjecture} \\
\text{generalized}\\\text{BG inequality}}
}
$$
And
\begin{align*}
Z_{\beta,H,\alpha,s}&=-\left(ch^B_3-\left(s+
\frac{1}{6}\right)\alpha^2
H^2ch_1^B)\right)\qquad s>0\\
&+i\left(Hch_2^B-\frac{\alpha^2}{2}
H^3ch_0\right)
\end{align*}

\medskip\noindent
{\bf GBG-inequality.} If $E$ is
$\nu_{B,H,\alpha}$-semistable,
\begin{equation}
\label{equation-GBG}
\begin{aligned}
H ch_2^B(E)-\frac{\alpha^2}{2}H^3
ch_0(E)&=0\\
\implies
ch_3^{B}-\frac{1}{6}\alpha^2H^2ch_1^B(E)&\leq
0
\end{aligned}
\end{equation}

\medskip\noindent
{\bf Bayer-Maci-Toda.} If Eq.
\ref{equation-GBG} is true then
$(\mathcal{A}^{B,H,\alpha},Z_{B,H,\alpha,s})$
is a strability condition.

\section{Remarks on
$\nu_{B,H,\alpha}$-stability}
\label{bsgt-label-004}
Keep in mind that
$$
\nu_{B,H,\alpha}=
\frac{Hch_2^B-\frac{\alpha^2}
{2}ch_0H^3}
{H^2 ch_1^B}
$$
and
$$
L \in \textit{Coh}^{B,H}(X) \qquad (L[1])
$$
as long as
$$
H^2ch_1^B(L)>0 \qquad (H^2ch_1^B(L[1])\leq 0.
$$
\medskip
If $E$ is $\nu_{B,H,\alpha}$-semistable for
$\alpha\gg 0$ then $E\in\textit{Coh}(X)$ and
is semistable with respect to the stability
given by
$$
q(t)=t \frac{H^2 ch_1}{ch_0}+ \frac{H
ch_2}{ch_0},\qquad t\gg 0
$$
and if $E$ is $q(t)$-semistable then $E$ is
$\nu_{B,H,\alpha}$-semistable for
$\frac{\alpha^2}{2}H^2>\mu_{B,H}(E)\Delta_{B,
H}(E)$.

If $X$ has Picard rank $\rho(X)=1$ then
$$
\Delta_{B,H}=H^3\Delta
$$
If $L$ is a line bundle
$\Delta_{B,H}(L)=H^3\Delta(L)=0$ (as $L[1])$,
then $L$ is $\nu_{B,H,\alpha}$-semistable as
long as $L \in \textit{Coh}^{B,H}(X)$ (as
$L[1])$

\section{Search for counterexamples}
\label{section-search-for-counterexamples}

{\bf Idea.} Look for counterexamples with
$ch_0=0$. When does $\mathcal{O}_D$ for
$D\geq 0$ divisor on $X$ (3-fold) produces a
counterexample?

Fix $H$ ample, $B=\beta H$. We need:
\begin{enumerate}
\item $\mathcal{O}_D$ is
$\nu_{\beta,\alpha}$-semistable.
\item $\nu_{\beta,\alpha}(\mathcal{O}_D)=0$.
\item
$ch_3^B(\mathcal{O}_D)-\frac{\alpha^2}{6}
H^2ch_1^B(\mathcal{O}_D)>0$.
\end{enumerate}

\begin{lemma}
\label{lemma-counterexample}
If $D \geq 0$ satisfies
$$
D^3>\frac{(DH^2)^3}{4(H^3)^2}+\frac{3}{4}
\frac{(D^3H)^2}{DH^2}
$$
then $\mathcal{O}_D$ is a counterexample for
the GBG inequality
\ref{equation-GBG}.
\end{lemma}

\begin{example}
\label{example-}
Let $Y$ be any smooth projective 3-fold,
$X=Bl_pY$, $D$ the exceptional divisor,
$L=\pi^*A$ with $A$ ample, $H=mL-D$ ample on
$X$.
\end{example}

[Slides]

\section{Concrete constructions}
\label{section-concrete-constructions}

How is a wall equation? Suppose that $\dim
X=2$, or even better that $X=
\mathbb{P}^2$. Also let
$\Gamma=\supset \oplus \supset \oplus
\frac{1}{2}\mathbb{Z}$ and
$\nu=(ch_0,ch_1,ch_2):=(r,c,d)$.

\begin{align*}
\nu_{\beta,\alpha}&
=\frac{ch_2^\beta-\frac{\alpha^2}{2}ch_0H^2}
{H
ch_1^\beta}\\
&=\frac{ch_2-\beta
ch_1+\left(\frac{\beta^2}{2}-
\frac{\alpha^2}{2}\right)ch_0}{ch_1-\beta
ch_0}
\end{align*}
A wall in the $(\beta,\alpha)$-plane is
produced by an exact sequence in
$\textit{Coh}^\beta(\mathbb{P}^2)$
$$
\xymatrix{
0\ar[r]&A\ar[r]&E\ar[r]&B\ar[r]&0
}
$$
with
$\nu_{\beta,\alpha}(A)=\nu_{\beta,\alpha}(E)$
and $ch(E)=V$.

\medskip\noindent
First: numerics: $w=(r',c',d')$ and solve
$\nu_{\beta,\alpha}(w)=\nu_{\beta,\alpha}(v)
$.

\begin{align*}
\beta^2-2\beta\left(\frac{rd'-r'd}{rc'-r'c}
\right)+\alpha^2&=
\frac{d c'-d'c}{rc'-r'c}
\end{align*}
$$
c=\frac{rd'-r'd}{rc'-r'c},\qquad
R=\sqrt{c^2+\frac{d c'-d'c}{rc'-r'c}}
$$
[Picture.] Look for wall for 1-dimensional
sheaves ($r=0$, $c>0$). $c=\frac{d}{c}$ same
center.

\begin{remark}
\label{bsgt-label-005}
Walls either coincide or they don't
intersect. [Picture.]
\end{remark}

\begin{lemma}[Bertan's lemma]
\label{lemma-Bertan}
If an object $A$ destabilizes $E$ at some
point then it destabilizes $E$ at every point
of the corresponding wall.
\end{lemma}

\medskip\noindent
Take $\beta=\frac{d}{c}$. If $A
\hookrightarrow E  \xrightarrow{surj.} B$ is
short exact in
$\textit{Coh}^\beta(\mathbb{P}^2)$.
\begin{itemize}
\item If $ch(E)=(0,c,d)$ and $E$ is
a sheaf then $A$ is a sheaf, just not a
subsheaf.
$$
\xymatrix{
0\ar[r]&\underbrace{\mathcal{H}^{-1}A}
_{\implies
0}\ar[r]
&  \underbrace{\mathcal{H}^{-1}E}_{=0}\ar[r]
&  \mathcal{H}^{-1}B\ar[r]
&  \mathcal{H}^0A\ar[r]&
\mathcal{H}^0E\ar[r] &
\mathcal{H}^0B\ar[r]&0
}
$$
$\implies$ $a= \mathcal{H}^0A$ sheaf.
\begin{enumerate}
\item $A$ produces a semi-circular wall only
when $ch_0(A)>0$.
\item $0 \leq  ch_1^\beta(A) \leq
ch_1^\beta(E)=c$.
This says that there are finitely many
possibilities for $c_1^\beta(A)$.
\item At the wall $W_{A,E}$ we have $A$ is
$\nu_{\beta,\alpha}$-semistable.
$$
\Delta_{\beta,\alpha}(A)\geq 0
$$
$$
ch_1^\beta(A)^2-2ch_0(A)ch_2^\beta(A)\geq 0
$$
$$
ch_1^\beta(A)^2\geq 2ch_0(A)ch_2^\beta(A)
$$
\item Wall equation
\begin{align*}
\nu_{\frac{1}{c}\alpha}(A)&=\nu_{\frac{1}{c},
\alpha}(E)\\
&=\frac{ch_2(E)-\beta ch_1(E)}{ch_1(E)}\\
&=\frac{d-\frac{d}{c}c}{c}=0
\end{align*}
\begin{align*}
ch_2^\beta(A)-\frac{\alpha^2}{2}ch_0^\beta(A)
&
=0\\
ch_2^\beta(A)=\frac{\alpha^2}{2}ch_2^\beta(A)
&
\geq 0
\end{align*}
$\implies$ {\bf walls are finite}.
\end{enumerate}
\end{itemize}

\medskip\noindent
Moduli of $\nu_{\beta,\alpha}$-semistable
objects are projective for $X=\mathbb{P}^2$.
[Picture of walls as nested circles in a
plane.] $\nu_{\beta \alpha}$-semistability
$\implies$ King's $\Theta$-semistability.

\begin{example}
\label{example-}
$v=(0,4,-4)$, $M_H(v)\leftarrow$ sheaves can
be found in one of the following short exact
sequences in
$\textit{Coh}^{-1}(\mathbb{P}^2)$.

[Duz\'et-Muican]

\noindent
Codimension 3:
$$
\xymatrix{
0\ar[r]&\mathcal{O}(1)\ar[r]&\mathcal{F}
\ar[r]
&\mathcal{O}(-2)\ar[r]&0
}
$$
Codimension 1:
$$
\xymatrix{
0\ar[r]&\mathcal{I}_P(1)\ar[r]&\mathcal{F}
\ar[r]
&\mathcal{I}_q^\vee(-3)[q]\ar[r]&0
}
$$
Codimension 0:
$$
\xymatrix{
0\ar[r]&2\mathcal{O}\ar[r]&\mathcal{F}\ar[r]&
\mathcal{O}(-2)[1]\ar[r]&0
}
$$
The first wall exchanges
$$
\xymatrix{
0\ar[r]&\mathcal{O}(1)\ar[r]&\mathcal{F}
\ar[r]
&\mathcal{O}(-3)[1]\ar[r]&0
}
$$
for (upon wall crossing)
$$
\xymatrix{
0\ar[r]&\mathcal{O}(-3)[1]\ar[r]&E^*\ar[r]&
\mathcal{O}(1)\ar[r]&0
}
$$
\begin{align*}
\text{Ext}^1(\mathcal{O}(1),\mathcal{O}(-3)
[1]
)
=\text{Ext}^2(\mathcal{O}(1),\mathcal{O}(-3))
\cong \Hom(\mathcal{O},\mathcal{O}(1))^\vee
\end{align*}
\end{example}

\begin{theorem}[Li-Zhao]
\label{theorem-Li-Zhao}
IF $v$ is primitive and $\sigma,\sigma'$ are
generic geometric with respect to $ v$ then
$\mathcal{M}_{\sigma}(v)$ and
$\mathcal{M}_{\sigma'}(v)$ are birational.
\end{theorem}

\medskip\noindent
We finish this minicourse with another
application of Schur's lemma:

\begin{theorem}[Koseki]
\label{theorem-Koseki}
Let $X$ be a smooth projective surface over
$k=\overline{k}$ with $\text{char}(k)>0$,
then there is a constant $C_{[X]}\geq 0$ and
dependending only on the birational
equivalence clas of $X$ such that
$$
\tilde{\Delta}(E)=\Delta(E)+C_{[X]}ch_0^2(E)
\geq
0
$$
for any $\mu_H$-semistable sheaf.
\end{theorem}

Morever
\begin{enumerate}
\item If $X$ is minimal and of general type,
$$
C_{[X]}=2+5K_X^2-\chi(\mathcal{O}_X)
$$
\item If $K(X)=1$ ($K$ is the Kodaira
dimension) and $X$ is quasi-elliptic then
$$
C_{[X]}=2-\chi(\mathcal{O}_X),\qquad
\text{[char 2+3]}
$$
\item $C_{[X]}=0$
\end{enumerate}

\medskip\noindent
{\bf Corollary.}
$$
Z_{B,H,\alpha}^{C_{[X]}}=-\left(\cong_2^B-
\left(\frac{C_{[X]}}{2H^2}+\frac{\alpha^2}{2}
\right)ch_0H^2\right)
+ich_1^B(X)
$$
is a stability condition on
$\textit{Coh}^{B,H}(X)$.

\begin{remark}
\label{remark-inequality}
$$
\tilde{\Delta}_{B,H}=\Delta_{B,H}+
ch_0^2C_{[X]
}\leq
2
$$
on any
$\nu_{B,H,\alpha}^{C_{[X]}}$-semistable
object.
\end{remark}

\begin{lemma}[Low key lemma]
\label{lemma-low-key-lemma}
\begin{itemize}
\item For $\alpha\gg 0$, $\nu$-semistable are
the Gieseker-semistable sheaves.
\item IF
$\alpha^2H^2>\mu_{B,H}(E)\tilde{\Delta}_{B,H}
(E)$
and if $E$ is
Gieseker-semistable $\implies$ $E$ is
Bridgeland semistable.
\end{itemize}
\end{lemma}

\begin{theorem}
\label{theorem-vanishing}
$H^2>6 C_{[X]}$ then $H^1(H \otimes K_X)=0$.
(If $\text{char}(k)=0$ we can take
$C_{[X]}=0$.)
\end{theorem}

\begin{proof}[Proof idea.]
Find a stability condition $\sigma=(Z,s)$
such that
\begin{enumerate}
\item $H$,  $\mathcal{O}[1] \in \mathcal{A}$
are $\sigma$-semistable.
\item $\mu_Z(H)>\mu_Z(\mathcal{O}[1])$
$$
\implies H^{1}(H \otimes K_X) \cong
\text{Ext}^1(H,\mathcal{O})^\vee
= \Hom(H_1\mathcal{O}[1])
$$
\end{enumerate}
\end{proof}
\end{document}
